\section*{Introduction}
Distributed acoustic sensing (DAS) converts an optical fibre into a dense array of strain or strain-rate sensors by interrogating Rayleigh backscatter from repeated laser pulses \citep{Zhan2020,LindseyMartin2021}. Because ordinary telecommunication cables can be used, DAS can provide metre-scale channel spacing over kilometres in places where conventional seismometers are sparse, such as cities \citep{Lindsey2017,Huot2024} and the seafloor \citep{Ide2021}. The price of this density is noise: channels are often poorly coupled to the ground, urban fibres record traffic and other anthropogenic signals, and individual channels are much noisier than inertial sensors \citep{Zhan2020,LindseyMartin2021}. Automatic event detection on DAS records therefore depends on exploiting the coherence of seismic arrivals across many channels.

Deep learning has become a standard tool for detection and phase picking on seismometer data \citep{ZhuBeroza2019,Mousavi2020,MousaviBeroza2022}, and several studies have transferred it to DAS. Convolutional neural networks (CNNs) that treat a space--time DAS window as an image have been used to detect local earthquakes on the Stanford telecom array \citep{Huot2024}, networks from the three classical deep-learning families, trained on seismometer records, have been tested on DAS measurements \citep{Hernandez2022}, and a phase picker for two-dimensional DAS data has been trained with semi-supervised learning \citep{Zhu2023}. Self-supervised networks have been proposed to suppress spatially incoherent DAS noise \citep{vandenEnde2023}.

DAS records are spatiotemporal fields, and at least four families of architectures encode this structure differently. A two-dimensional CNN applies local filters jointly in space and time. A CNN--LSTM uses a CNN to summarise the channels at each time step and a long short-term memory (LSTM) network \citep{Hochreiter1997} to model the temporal sequence. A convolutional LSTM (ConvLSTM) keeps the spatial layout inside its recurrent state \citep{Shi2015}. Transformer encoders \citep{Vaswani2017}, applied to space--time patches as in the Vision Transformer \citep{Dosovitskiy2021}, learn interactions between all patches through self-attention. Many published DAS studies present a single architecture. Differences in data, preprocessing, training budget and evaluation then make it hard to know whether a reported gain comes from the architecture.

This report compares the four families under a single protocol on a public, labelled DAS earthquake-detection dataset recorded on telecommunication fibre under the Stanford University campus \citep{HuotData2022,Huot2024}. We ask three practical questions: (1) with all labels and the same training budget, do the architectures differ in detection skill and computational cost; (2) how does each degrade when synthetic noise or dead channels are added to the real test records; and (3) how does each behave when only 1\,\% or 10\,\% of the training labels are available? All four models are deliberately small and were trained on a CPU, a setting typical of student and small-observatory computing resources. We report seed-to-seed variability and paired statistical tests, and we state clearly what the dataset cannot answer.

\section*{Data}
\subsection*{Stanford telecom-fibre DAS earthquake-detection dataset}
We use the dataset ``Stanford fiber-optic DAS array: Earthquake detection dataset'' \citep{HuotData2022}, distributed under a CC BY 4.0 licence on Zenodo (\url{https://doi.org/10.5281/zenodo.6344570}). The data are real field recordings; no synthetic waveforms are used anywhere in this study except the noise perturbations described below. The recordings were made from 2016 to 2019 with a DAS interrogator connected to fibres in telecommunication conduits under the Stanford University campus \citep{Huot2024}. According to the preprint describing the dataset \citep{HuotArxiv2022}, the array used an OptaSense ODH3 interrogator, channel spacing of 8\,m, and an approximately 2.4-km fibre path laid out in a figure-of-eight; data were recorded at 50 samples\,s$^{-1}$. Earthquake labels were derived from the Northern California Earthquake Data Center (NCEDC) catalogue for events within 100\,km, and noise windows were drawn at times at least 5 minutes away from catalogued events and checked visually for earthquake signals \citep{HuotArxiv2022}.

The distributed files are TensorFlow Record shards already split into training (10 shards), evaluation (1 shard) and test (1 shard) sets. We use these splits unchanged. The companion code \citep{HuotCode2022} documents the preprocessing that produced each example: conversion to strain rate, removal of the median across channels at each time sample, 1--12\,Hz band-pass filtering, decimation by a factor of 2 (to 25 samples\,s$^{-1}$), cropping to a fixed window around the predicted arrival time, extraction of two 288-channel fibre segments (channels 13--300 and 328--615, the second one reversed) from each window, and clipping at $\pm0.0299$ followed by division by 0.01572. We did not repeat these steps; we read the distributed values directly.

After downloading, we verified every file against the MD5 checksum published on Zenodo. Each example is a 288-channel $\times$ 695-sample float32 array (27.8\,s at 25 samples\,s$^{-1}$) with a binary label (1 = earthquake, 0 = noise). The training, evaluation and test sets contain %%NTRAIN%%, %%NEVAL%% and %%NTEST%% examples respectively, each exactly balanced between the two classes (Table~\ref{tab:data}). All values lie in $[-1.90, 1.90]$, there are no missing values, and a hash comparison found no identical examples within or across the splits. Figure~\ref{fig:data} shows one earthquake and one noise example from the test set.

\subsection*{Model input}
To keep training feasible on a CPU, every example is reduced by non-overlapping average pooling over 6 adjacent channels (48\,m of fibre) and 2 adjacent samples, giving 48 channels $\times$ 347 samples at 12.5 samples\,s$^{-1}$. Averaging adjacent samples acts as a crude low-pass filter, so energy above about 6\,Hz is attenuated; the spatial averaging suppresses incoherent channel noise while preserving arrivals whose apparent wavelength along the fibre is much longer than 48\,m, which is the case for regional earthquakes recorded on a horizontal fibre. The same reduced input is given to all four models. During training a 320-sample (25.6\,s) window is cropped at a random position and its polarity is flipped with probability 0.5; at evaluation the central 320 samples are used.

\section*{Methods}
\subsection*{Architectures}
All models take a $1\times48\times320$ space--time patch and output one logit for the probability that the patch contains an earthquake (Table~\ref{tab:arch}). Sizes were chosen to fall within the same order of magnitude of parameters (46--150 thousand) and of forward cost (0.05--0.12 GFLOPs per patch).
\begin{itemize}
\item \textbf{CNN.} Four blocks of $3\times3$ convolution, batch normalisation \citep{Ioffe2015}, ReLU and $2\times2$ max pooling (16, 32, 48 and 64 filters), global average pooling, dropout (0.3) and a linear output.
\item \textbf{CNN--LSTM.} Three such convolutional blocks (16, 32, 48 filters) reduce the patch to $48\times6\times40$; features are averaged over the spatial axis and the 40-step sequence is processed by a bidirectional LSTM with 64 units per direction, whose outputs are averaged over time before dropout and a linear output.
\item \textbf{ConvLSTM.} Two convolutional blocks (16, 32 filters) reduce the patch to $32\times12\times80$. A ConvLSTM cell \citep{Shi2015} with 32 hidden channels and one-dimensional convolutions (kernel 5) along the fibre then runs over the 80 time steps, so the recurrent state keeps its spatial layout. Hidden states are averaged over space and time before the output layer.
\item \textbf{Transformer.} The patch is split into non-overlapping $8\times16$ (channel $\times$ sample) patches, each linearly embedded into 64 dimensions, giving 120 tokens plus a learnable class token and learnable position embeddings \citep{Dosovitskiy2021}. Four pre-normalisation encoder layers \citep{Vaswani2017} with 4 heads, feed-forward width 128, GELU activation and dropout 0.1 are followed by layer normalisation and a linear output on the class token.
\end{itemize}

\subsection*{Training protocol}
All models were trained with the same protocol and budget: binary cross-entropy loss, AdamW \citep{Loshchilov2019} with weight decay $10^{-2}$, a one-cycle learning-rate schedule \citep{SmithTopin2019} with maximum learning rate $10^{-3}$ and 10\,\% warm-up, batch size 32, gradient-norm clipping at 1, and 20 epochs (%%STEPS%% optimiser steps for the full training set). For the reduced-label experiments the number of steps was raised to a minimum of 300 so that small training sets still receive a usable number of updates. The model was evaluated on the evaluation split once per epoch (at least every 10 steps), and the weights with the highest evaluation ROC-AUC were kept; the test split was used only once, after training. Each configuration was trained with five random seeds (0--4) that control weight initialisation, mini-batch order, cropping, polarity flips and, for reduced labels, which examples are kept. Every run used one CPU thread of the same virtual machine (Intel Xeon 2.1\,GHz, 2 virtual cores, 7\,GB RAM, no GPU), two runs in parallel. Models were implemented in PyTorch \citep{Paszke2019}; metrics and tests used scikit-learn \citep{Pedregosa2011} and SciPy \citep{Virtanen2020}.

\subsection*{Experiments}
\textbf{E1, full labels.} Each architecture is trained on all %%NTRAIN%% training examples. We report test ROC-AUC, accuracy and F1 score at a 0.5 threshold, together with cost: trainable parameters, forward floating-point operations per patch (counted with PyTorch's FLOP counter), median single-patch inference latency, median time of one training step with batch 32, increase of peak resident memory during training, and wall-clock training time.

\textbf{E2, noise robustness.} The five E1 models of each architecture are evaluated, without retraining, on perturbed copies of the real test set. Perturbations are synthetic and are applied at full resolution before pooling. (a) Additive zero-mean white Gaussian noise at signal-to-noise ratios of 20, 10, 5, 0 and $-5$\,dB, where SNR $=10\log_{10}(P_\mathrm{patch}/P_\mathrm{noise})$ and $P_\mathrm{patch}$ is the mean square of the original patch (which itself contains the recorded background noise). (b) Channel dropout, in which 10, 30 or 50\,\% of the 288 channels, chosen at random, are set to zero to mimic dead or decoupled channels. The same random perturbation is used for all models, so comparisons are paired.

\textbf{E3, label scarcity.} Each architecture is retrained from scratch on stratified random subsets containing 1\,\% (%%N001%% examples) and 10\,\% (%%N01%% examples) of the training set, with the evaluation and test sets unchanged.

\subsection*{Statistical analysis}
The test set contains %%NTEST%% examples, so differences of a few percentage points are within sampling noise. We therefore use three complementary analyses. Across seeds, we report mean $\pm$ standard deviation and apply the Friedman rank test \citep{Friedman1937}, treating seeds as blocks, as recommended for comparing several classifiers \citep{Demsar2006}. On the test examples, we average the five seeds' probabilities for each architecture (a seed ensemble) and compare architectures pairwise with (i) a stratified paired bootstrap of the difference in ROC-AUC (5000 resamples; percentile 95\,\% interval; two-sided $p$-value) \citep{EfronTibshirani1993} and (ii) the exact McNemar test on the 0.5-threshold decisions \citep{McNemar1947}. $p$-values over the six pairs are adjusted with the Holm procedure \citep{Holm1979}. We consider adjusted $p<0.05$ as evidence of a difference.
